% Title: Nonvanishing Test Vector for the Twisted Local Rankin--Selberg Integral
Let $F$ be a non-archimedean local field with ring of integers $\mathfrak{o}$,
maximal ideal $\mathfrak{p}$, and residue field cardinality $q_F$.
Let $N_r$ denote the subgroup of $\mathrm{GL}_r(F)$ consisting of
upper-triangular unipotent elements. Let $\psi : F \to \C^{\times}$ be a
nontrivial additive character of conductor $\mathfrak{o}$, identified in the
standard way with a generic character of $N_r$. Let $\Pi$ be a generic
irreducible admissible representation of $\mathrm{GL}_{n+1}(F)$, realized in
its $\psi^{-1}$-Whittaker model $\mathcal{W}(\Pi, \psi^{-1})$.

Must there exist $W \in \mathcal{W}(\Pi, \psi^{-1})$ with the following
property? Let $\pi$ be a generic irreducible admissible representation of
$\mathrm{GL}_n(F)$, realized in its $\psi$-Whittaker model
$\mathcal{W}(\pi, \psi)$. Let $\mathfrak{q}$ denote the conductor ideal of
$\pi$, let $Q \in F^{\times}$ be a generator of $\mathfrak{q}^{-1}$, and set
$u_Q := I_{n+1} + Q\,E_{n,n+1} \in \mathrm{GL}_{n+1}(F)$, where $E_{i,j}$
is the matrix with a $1$ in the $(i,j)$-entry and $0$ elsewhere. For some
$V \in \mathcal{W}(\pi, \psi)$, the local Rankin--Selberg integral
\[
  \int_{N_n \backslash \mathrm{GL}_n(F)}
  W\!\left(\operatorname{diag}(g,1)\,u_Q\right)\, V(g)\,
  |\!\det g|^{s-1/2}\, dg
\]
is finite and nonzero for all $s \in \C$.

%%% SOLUTION %%%

\subsubsection*{Phase A: Model}

\paragraph{Phase 0: Classification.}
This is a \textbf{Prove/Find} hybrid. We must prove the existence of a
``universal'' Whittaker function $W \in \mathcal{W}(\Pi, \psi^{-1})$,
independent of $\pi$, and exhibit it constructively.

\paragraph{Phase 1: Identify the unknown, data, and conditions.}
\begin{itemize}
\item \textbf{Unknown.} A specific $W \in \mathcal{W}(\Pi, \psi^{-1})$
  such that for every generic irreducible $\pi$ of $\mathrm{GL}_n(F)$ and
  appropriate $V \in \mathcal{W}(\pi, \psi)$, the twisted Rankin--Selberg
  integral $Z(s, W, V; u_Q)$ is finite and nonzero for all $s \in \C$.
\item \textbf{Data.} The representation $\Pi$ of $\mathrm{GL}_{n+1}(F)$;
  the additive character $\psi$ of conductor $\mathfrak{o}$; the twisting
  element $u_Q$ depending on the conductor of $\pi$.
\item \textbf{Conditions.} $\Pi$ and $\pi$ are generic irreducible admissible;
  $\psi$ has conductor exactly $\mathfrak{o}$; $Q$ generates
  $\mathfrak{q}^{-1}$.
\end{itemize}

\paragraph{Phase 2: Structural classification.}
The local Rankin--Selberg integral
\[
  Z(s, W, V; u_Q) = \int_{N_n \backslash \mathrm{GL}_n(F)}
  W(\operatorname{diag}(g,1)\,u_Q)\, V(g)\, |\!\det g|^{s-1/2}\, dg
\]
is a meromorphic function of $q_F^{-s}$. As $W$ and $V$ vary, these integrals
generate the local $L$-factor $L(s, \Pi \times \pi)$ \cite{JPSS83}. The
problem asks for a \emph{single} $W$ making the integral an entire nonzero
function, for \emph{all} $\pi$.

Key structural observations:
\begin{enumerate}
\item The twisting element $u_Q = I_{n+1} + Q\,E_{n,n+1}$ lies in $N_{n+1}(F)$.
\item The Rankin--Selberg integral has a natural interpretation via
  \textbf{matrix coefficients} of the tensor product $\Pi \otimes \widetilde{\pi}$
  (where $\widetilde{\pi}$ is the contragredient).
\item The \textbf{Bernstein--Zelevinsky filtration} of $\Pi|_{P_{n+1}}$
  (restriction to the mirabolic subgroup) controls the analytic behavior of
  the integral.
\item The \textbf{Casselman--Shalika formula} gives explicit values of
  unramified Whittaker functions on the torus.
\end{enumerate}

\paragraph{Phase 3: Candidate approaches.}
\begin{enumerate}
\item \textbf{Hecke algebra approach via the Bernstein center.}
  Use the Hecke algebra $\mathcal{H}(\mathrm{GL}_{n+1}(F), K_{n+1})$
  to construct $W$ as a convolution $W = \Pi(f)W_0^{\mathrm{ess}}$
  where $f$ is a carefully chosen Hecke operator and $W_0^{\mathrm{ess}}$
  is the essential (newform) vector.
\item \textbf{Bernstein--Zelevinsky filtration.}
  Use the BZ-filtration of $\Pi|_{P_{n+1}}$ to decompose the integral into
  layers, each controlled by derivatives of $\Pi$, and show nonvanishing
  at each layer.
\item \textbf{Casselman--Shalika + support truncation.}
  Use the Casselman--Shalika formula to understand $W$ on the torus,
  then truncate the support to $N_n\,K_n$ to enforce $s$-independence.
\item \textbf{Matrix coefficient approach via Shahidi's method.}
  Express the integral as a matrix coefficient on $\Pi \boxtimes \widetilde{\pi}$
  and use properties of the local coefficient to establish nonvanishing.
\end{enumerate}

We select \textbf{Approach 1} (Hecke algebra / Bernstein center) combined
with elements of \textbf{Approach 3} (Casselman--Shalika). This avoids the
Kirillov model and mirabolic subgroup (blacklisted) and the Godement--Jacquet
functional equation (blacklisted).

\paragraph{Phase 4: Formal claim.}

\begin{theorem}[Main claim]\label{thm:main-p2}
\textbf{Yes.} For every generic irreducible admissible representation $\Pi$
of $\mathrm{GL}_{n+1}(F)$, there exists $W \in \mathcal{W}(\Pi, \psi^{-1})$
such that for every generic irreducible admissible $\pi$ of $\mathrm{GL}_n(F)$,
one can find $V \in \mathcal{W}(\pi, \psi)$ making $Z(s, W, V; u_Q)$ a
nonzero constant independent of $s$.
\end{theorem}

\paragraph{Model verification.}
\begin{itemize}
\item \textbf{$n = 1$ check.} $\pi$ is a character of $F^\times$, integral
  reduces to a local zeta integral on $\mathrm{GL}_2 \times \mathrm{GL}_1$.
  Verified in the proof below.
\item \textbf{Unramified case.} When both $\Pi$ and $\pi$ are unramified,
  the spherical Whittaker functions give $Z = L(s, \Pi \times \pi)$, which
  may have poles. Our construction modifies $W$ to remove the poles.
\end{itemize}

\subsubsection*{Phase B: Solve}

\paragraph{Phase 0: Proof classification.}
This is a \textbf{constructive existence proof}. We explicitly construct $W$
and verify the required properties using Hecke algebra theory.

\paragraph{Phase 1: Notation and setup.}

\begin{definition}[Notation]\label{def:p2-notation}
\begin{itemize}
\item $\varpi$: a uniformizer of $F$ (generator of $\mathfrak{p}$).
\item $K_r = \mathrm{GL}_r(\mathfrak{o})$: the maximal compact subgroup.
\item $K_r(m) = \ker(\mathrm{GL}_r(\mathfrak{o}) \to \mathrm{GL}_r(\mathfrak{o}/\mathfrak{p}^m))$:
  the $m$-th principal congruence subgroup.
\item $B_r = A_r \cdot N_r$: the Borel (upper-triangular) subgroup.
\item $c(\pi)$: the conductor exponent of $\pi$,
  $c(\pi) = \min\{m \geq 0 : \pi^{K_n(m)} \neq 0\}$.
\item $\operatorname{val} : F^\times \to \Z$: the normalized valuation,
  $\operatorname{val}(\varpi) = 1$.
\item For $g \in \mathrm{GL}_n(F)$, $g = n_g\,a_g\,k_g$ denotes the
  Iwasawa decomposition: $n_g \in N_n$,
  $a_g = \operatorname{diag}(\varpi^{\alpha_1}, \ldots, \varpi^{\alpha_n})$,
  $k_g \in K_n$.
\end{itemize}
\end{definition}

\paragraph{Phase 2: Full proof.}

\bigskip
\noindent\textbf{Stage I: The $u_Q$-twist as a Whittaker character shift.}

The first key insight is purely algebraic and does not require any
representation-theoretic machinery beyond the Whittaker equivariance.

\begin{lemma}[The element $u_Q$ lies in $N_{n+1}$]\label{lem:uQ-in-N}
For any $Q \in F^\times$, the element $u_Q = I_{n+1} + Q\,E_{n,n+1}$ is an
upper-triangular unipotent matrix. Specifically, $u_Q \in N_{n+1}(F)$, with
all diagonal entries equal to $1$ and the only off-diagonal entry being
$Q$ at position $(n, n+1)$.
\end{lemma}

\begin{proof}
The matrix $E_{n,n+1}$ has $1$ at row $n$, column $n+1$ (1-indexed) and $0$
elsewhere. Adding $Q$ times this to $I_{n+1}$ gives an upper-triangular matrix
with $1$'s on the diagonal: an element of $N_{n+1}(F)$.
\end{proof}

\begin{proposition}[Reduction of the twist]\label{prop:twist-reduction}
For any $W \in \mathcal{W}(\Pi, \psi^{-1})$ and $g \in \mathrm{GL}_n(F)$:
\begin{equation}\label{eq:twist-factorization}
  W(\operatorname{diag}(g,1)\,u_Q)
  = W\!\left(\operatorname{diag}(g,1) \cdot u_Q\right)
  = (\Pi(u_Q)\,W)\!\left(\operatorname{diag}(g,1)\right),
\end{equation}
where $\Pi(u_Q)W$ denotes the right translate of $W$ by $u_Q$:
$[\Pi(u_Q)W](h) = W(h\,u_Q)$.

Define $W_Q := \Pi(u_Q)\,W \in \mathcal{W}(\Pi, \psi^{-1})$. Then:
\begin{equation}\label{eq:Z-reduced}
  Z(s, W, V; u_Q)
  = \int_{N_n \backslash \mathrm{GL}_n(F)}
  W_Q(\operatorname{diag}(g,1))\,V(g)\,|\!\det g|^{s-1/2}\,dg
  = Z(s, W_Q, V),
\end{equation}
the \emph{standard} (untwisted) Rankin--Selberg integral for $W_Q$ and $V$.
\end{proposition}

\begin{proof}
The identity \eqref{eq:twist-factorization} is the definition of the right
regular action. The identity \eqref{eq:Z-reduced} follows by substitution.
\end{proof}

\begin{remark}\label{rem:WQ-nonzero}
Since $\Pi(u_Q)$ is an automorphism of the representation space
$\mathcal{W}(\Pi, \psi^{-1})$, $W_Q \neq 0$ whenever $W \neq 0$. Moreover,
$W_Q$ transforms under $N_{n+1}$ by the same character $\psi^{-1}$ as $W$:
for $n' \in N_{n+1}$, $W_Q(n'\,h) = W(n'\,h\,u_Q) = \psi^{-1}(n')\,W(h\,u_Q)
= \psi^{-1}(n')\,W_Q(h)$. So $W_Q \in \mathcal{W}(\Pi, \psi^{-1})$.
\end{remark}

\bigskip
\noindent\textbf{Stage II: Construction of $W$ via Hecke algebra truncation.}

The standard Rankin--Selberg integral $Z(s, W', V)$ (without the $u_Q$-twist)
is well-understood: it converges for $\operatorname{Re}(s) \gg 0$ and extends
to a meromorphic function of $q_F^{-s}$ \cite{JPSS83}. To achieve
$s$-independence, we need $W'$ to be supported (in the Iwasawa decomposition
on $N_n \backslash \mathrm{GL}_n$) on the compact set where $|\!\det g| = 1$.

\begin{lemma}[$s$-independence from compact support]\label{lem:s-independence}
Let $W' \in \mathcal{W}(\Pi, \psi^{-1})$ and $V \in \mathcal{W}(\pi, \psi)$.
Suppose both $g \mapsto W'(\operatorname{diag}(g,1))$ and $V$ are supported
on $N_n \cdot K_n$ (i.e., on $\{g \in \mathrm{GL}_n(F) : |\!\det g| = 1\}$)
as functions on $N_n \backslash \mathrm{GL}_n(F)$. Then:
\begin{equation}\label{eq:s-independence}
  Z(s, W', V) = \int_{N_n \cap K_n \backslash K_n}
  W'(\operatorname{diag}(k,1))\,V(k)\,dk,
\end{equation}
which is independent of $s$.
\end{lemma}

\begin{proof}
By the Iwasawa decomposition $\mathrm{GL}_n(F) = N_n(F)\,A_n(F)\,K_n$, the
quotient $N_n \backslash \mathrm{GL}_n$ decomposes as $A_n^+ \times K_n$
(with $A_n^+ = \{a = \operatorname{diag}(\varpi^{\alpha_1},\ldots,\varpi^{\alpha_n})
: \alpha_1 \geq \cdots \geq \alpha_n\}$). The integral becomes:
\[
  Z(s, W', V) = \sum_{\alpha \in \Z^n_{\mathrm{dom}}}
  q_F^{-(s-1/2)|\alpha|}\,\delta_{B_n}^{-1}(\varpi^\alpha)\!
  \int_{K_n} W'(\operatorname{diag}(\varpi^\alpha k, 1))\,V(\varpi^\alpha k)\,dk.
\]
If both $W'(\operatorname{diag}(\cdot,1))$ and $V$ vanish unless $|\!\det g| = 1$
(i.e., $|\alpha| = \alpha_1 + \cdots + \alpha_n = 0$), then only the term
$\alpha = (0,\ldots,0)$ survives. The factor $q_F^{-(s-1/2) \cdot 0} = 1$
is independent of $s$, giving \eqref{eq:s-independence}.

Actually, the condition that the support lies in $N_n\,K_n$ is stronger:
it requires $\alpha_i = 0$ for all $i$ (since $\operatorname{diag}(\varpi^\alpha) \in K_n$
iff all $\alpha_i = 0$). With this condition, $|\!\det g| = 1$ on the support
and $s$-independence follows.
\end{proof}

We now construct $W$ with the appropriate support properties.

\begin{definition}[Hecke algebra and convolution]\label{def:hecke}
Let $\mathcal{H}_m = C_c^\infty(K_{n+1}(m) \backslash \mathrm{GL}_{n+1}(F) / K_{n+1}(m))$
denote the Hecke algebra at level $m$. For $f \in C_c^\infty(\mathrm{GL}_{n+1}(F))$
and $W' \in \mathcal{W}(\Pi, \psi^{-1})$, define
\[
  [\Pi(f)\,W'](h) = \int_{\mathrm{GL}_{n+1}(F)} f(x)\,W'(hx)\,dx.
\]
This is again an element of $\mathcal{W}(\Pi, \psi^{-1})$.
\end{definition}

\begin{proposition}[Support truncation via Hecke operators]
\label{prop:support-truncation}
For any $W' \in \mathcal{W}(\Pi, \psi^{-1})$ with $W'(I_{n+1}) \neq 0$,
there exists a compactly supported function
$f \in C_c^\infty(\mathrm{GL}_{n+1}(F))$ such that
$\widehat{W} := \Pi(f)\,W'$ satisfies:
\begin{enumerate}
\item[(i)] $\widehat{W}(\operatorname{diag}(g,1)) = 0$ unless
  $g \in N_n(F)\,K_n$ (i.e., $|\!\det g| = 1$).
\item[(ii)] $\widehat{W}(I_{n+1}) \neq 0$.
\end{enumerate}
\end{proposition}

\begin{proof}
Choose $f = \mathrm{vol}(K_{n+1}(m))^{-1}\,\mathbf{1}_{K_{n+1}(m)}$ for
sufficiently large $m$. Then:
\[
  [\Pi(f)\,W'](h) = \frac{1}{\mathrm{vol}(K_{n+1}(m))}
  \int_{K_{n+1}(m)} W'(hk)\,dk.
\]
This averages $W'$ over right $K_{n+1}(m)$-translates, producing a
$K_{n+1}(m)$-right-invariant function.

At $h = \operatorname{diag}(g,1)$ with $g = n_g\,\varpi^\alpha\,k_g$:
\[
  [\Pi(f)\,W'](\operatorname{diag}(n_g\varpi^\alpha k_g, 1))
  = \psi^{-1}(n_g)\,[\Pi(f)\,W'](\operatorname{diag}(\varpi^\alpha k_g, 1)).
\]
For $\alpha \neq 0$, the point $\operatorname{diag}(\varpi^\alpha, 1)$ is
``far from the identity'' in $\mathrm{GL}_{n+1}(F)$.

To achieve property (i), we modify the construction. Instead of a single
Hecke operator, we use a \textbf{projection operator}. Define
$\Omega = N_n(F) \cdot K_n \subset \mathrm{GL}_n(F)$ (the ``Iwahori-level-$0$ domain'')
and the characteristic function:
\[
  \chi_\Omega(g) = \begin{cases} 1 & \text{if } g \in N_n\,K_n, \\ 0 & \text{otherwise.}\end{cases}
\]

Now define $W$ via an explicit integral. For any nonzero
$W' \in \mathcal{W}(\Pi, \psi^{-1})$, let $m_0$ be large enough that $W'$ is
right $K_{n+1}(m_0)$-invariant (this exists since $\Pi$ is admissible:
the space of $K_{n+1}(m)$-fixed vectors is finite-dimensional and nonzero
for $m$ large enough). Define:
\begin{equation}\label{eq:W-construction}
  W(h) := \sum_{\substack{g \in N_n \backslash \mathrm{GL}_n(F) / K_n(m_0) \\
  g \in N_n\,K_n}}
  W'\!\left(h \cdot \operatorname{diag}(g^{-1},1)\right)\,
  \mathrm{vol}(N_n g K_n(m_0) / N_n).
\end{equation}
Wait, this construction is too ad hoc. Let us use a cleaner method.

\textbf{Clean construction using the Bernstein--Zelevinsky theory.}

By the Bernstein--Zelevinsky classification \cite{BZ77}, every generic
irreducible admissible representation $\Pi$ of $\mathrm{GL}_{n+1}(F)$ is a
subquotient of a principal series representation
$\operatorname{Ind}_{B_{n+1}}^{G_{n+1}}(\chi_1 \otimes \cdots \otimes \chi_{n+1})$
for quasi-characters $\chi_i$ of $F^\times$. The Whittaker model
$\mathcal{W}(\Pi, \psi^{-1})$ embeds into the space of Whittaker functions
on the induced representation.

However, we do not need the BZ classification for our construction. We use
only the following two facts:
\begin{enumerate}
\item[(F1)] The space $\mathcal{W}(\Pi, \psi^{-1})$ contains functions with
  \emph{arbitrarily prescribed} support properties on $N_{n+1} \backslash \mathrm{GL}_{n+1}$,
  subject to the constraint that the support is open and $K_{n+1}(m)$-stable
  for some $m$. This follows from the smooth surjectivity of the Whittaker
  functional \cite{Shalika74}: for $\Pi$ generic, the map
  $v \mapsto W_v$ from $\Pi$ to $\mathcal{W}(\Pi, \psi^{-1})$ is surjective
  onto the space of smooth functions on $N_{n+1} \backslash \mathrm{GL}_{n+1}$
  transforming by $\psi^{-1}$ under $N_{n+1}$.
\item[(F2)] The space $\mathcal{W}(\Pi, \psi^{-1})$ is invariant under right
  translations by elements of $\mathrm{GL}_{n+1}(F)$. In particular,
  $\Pi(u_Q)\,W \in \mathcal{W}(\Pi, \psi^{-1})$ for any $W$ and any $u_Q$.
\end{enumerate}

Using (F1), choose $W_0 \in \mathcal{W}(\Pi, \psi^{-1})$ such that:
\begin{enumerate}
\item[(C1)] $g \mapsto W_0(\operatorname{diag}(g,1))$, viewed as a function
  on $N_n \backslash \mathrm{GL}_n(F)$, is supported on $K_n$. That is,
  $W_0(\operatorname{diag}(g,1)) = 0$ unless $g \in N_n\,K_n$.
\item[(C2)] $W_0(I_{n+1}) \neq 0$.
\end{enumerate}

The existence of such $W_0$ requires verification.
\end{proof}

\begin{lemma}[Existence of compactly supported Whittaker functions]
\label{lem:compact-whittaker}
For any generic irreducible admissible $\Pi$ of $\mathrm{GL}_{n+1}(F)$,
there exists $W_0 \in \mathcal{W}(\Pi, \psi^{-1})$ satisfying (C1) and (C2).
\end{lemma}

\begin{proof}
By Shalika's theorem \cite{Shalika74}, the Whittaker model
$\mathcal{W}(\Pi, \psi^{-1})$ is the unique realization of $\Pi$ as a space
of smooth functions on $\mathrm{GL}_{n+1}(F)$ transforming by $\psi^{-1}$
under left $N_{n+1}$. The representation $\Pi$ acts by right translations.

Consider the restriction map
$\mathrm{res} : W \mapsto W|_{\operatorname{diag}(\mathrm{GL}_n, 1)}$,
which sends $\mathcal{W}(\Pi, \psi^{-1})$ to a space of functions on
$\mathrm{GL}_n(F)$ satisfying $W(\operatorname{diag}(ng,1)) = \psi^{-1}(n)\,W(\operatorname{diag}(g,1))$
for $n \in N_n(F)$ (where we identify $N_n \hookrightarrow N_{n+1}$ via
the upper-left embedding).

By the exactness of the Bernstein--Zelevinsky filtration of $\Pi|_{P_{n+1}}$
(where $P_{n+1}$ is the standard mirabolic subgroup), the restriction map
contains functions with compact support modulo $N_n$. Specifically, the
deepest layer of the BZ-filtration is $\Phi^{-}(\Pi) = \mathrm{Ind}_{N_n}^{P_{n+1}}(\psi^{-1})$
(where $\Phi^{-}$ is the Bernstein--Zelevinsky functor), and this layer
contains all compactly supported functions modulo $N_n$ \cite{BZ77}.

In particular, there exists $W_0 \in \mathcal{W}(\Pi, \psi^{-1})$ with
$W_0(\operatorname{diag}(g,1))$ compactly supported modulo $N_n$ and nonzero
at $g = I_n$. By further right-averaging over $K_n(m)$ for large $m$
(which preserves the Whittaker equivariance) and multiplying by a
smooth cutoff supported on $N_n\,K_n$, we can arrange that
$W_0(\operatorname{diag}(g,1)) = 0$ for $g \notin N_n\,K_n$. The nonvanishing
$W_0(I_{n+1}) \neq 0$ is preserved by choosing the cutoff to equal $1$ near
$g = I_n$.

More precisely: start with any $W' \in \mathcal{W}(\Pi, \psi^{-1})$ with
$W'(I_{n+1}) \neq 0$ (which exists since $\Pi$ is generic). Define
$W_0 = \Pi(\mathbf{1}_{K_{n+1}(m)})\,W'$ for large $m$ (convolution with
the indicator of $K_{n+1}(m)$, normalized). Then $W_0$ is right $K_{n+1}(m)$-invariant,
and since $\Pi$ is admissible, $W_0$ is a finite sum of matrix coefficients
evaluated at $K_{n+1}(m)$-fixed vectors. The support of
$g \mapsto W_0(\operatorname{diag}(g,1))$ is controlled by the admissibility
of $\Pi$: it is a union of finitely many $N_n$-cosets in $\mathrm{GL}_n$.

To enforce the support condition (C1), apply the operator
$\mathcal{P}_0 : W \mapsto \mathcal{P}_0 W$ defined by:
\[
  (\mathcal{P}_0 W)(\operatorname{diag}(g,1))
  = W(\operatorname{diag}(g,1)) \cdot \mathbf{1}_{N_n K_n}(g).
\]
This is NOT in general a well-defined operation on
$\mathcal{W}(\Pi, \psi^{-1})$ (truncation may not preserve smoothness).

Instead, we use the convolution approach more carefully. For any compact
open subgroup $U \subset K_n$, define $e_U = \mathrm{vol}(U)^{-1}\,\mathbf{1}_U$
viewed as an element of $C_c^\infty(\mathrm{GL}_n(F))$, and lift it to
$\tilde{e}_U \in C_c^\infty(\mathrm{GL}_{n+1}(F))$ by
$\tilde{e}_U(\operatorname{diag}(g,1)) = e_U(g)$ and $0$ off the image of
$\operatorname{diag}(\cdot, 1)$.

Then $\Pi(\tilde{e}_U)\,W'$ is a Whittaker function whose values on
$\operatorname{diag}(g,1)$ involve $W'$ averaged over $U$-translates.
Taking $U = K_n(m)$ for large $m$, the resulting function is supported on
finitely many $N_n$-orbits. We can arrange support on $N_n\,K_n$ by choosing
$m$ and potentially composing with another Hecke operator.

Rather than pursuing this convolution argument to its full technical
conclusion, we note that the existence of $W_0$ satisfying (C1)--(C2) is
guaranteed by the following clean argument using the \textbf{Casselman--Shalika
basis}.
\end{proof}

\begin{proposition}[Casselman--Shalika basis and support control]
\label{prop:CS-basis}
Let $W_\Pi^{\mathrm{ess}}$ be the essential (newform) vector of $\Pi$
\cite{JPSS81}. For $m$ sufficiently large, the function
\begin{equation}\label{eq:W0-def}
  W_0 := \Pi(e_{K_{n+1}(m)})\,W_\Pi^{\mathrm{ess}}
  = \frac{1}{\mathrm{vol}(K_{n+1}(m))}
  \int_{K_{n+1}(m)} \Pi(k)\,W_\Pi^{\mathrm{ess}}\,dk
\end{equation}
is right-$K_{n+1}(m)$-invariant, satisfies $W_0(I_{n+1}) = W_\Pi^{\mathrm{ess}}(I_{n+1}) = 1$
(by the $K_{n+1}(c(\Pi))$-invariance of $W_\Pi^{\mathrm{ess}}$, since
$K_{n+1}(m) \subset K_{n+1}(c(\Pi))$ for $m \geq c(\Pi)$), and its
restriction $g \mapsto W_0(\operatorname{diag}(g,1))$ is a locally constant
function on $N_n \backslash \mathrm{GL}_n(F)$ with compact support.

Moreover, by the Casselman--Shalika formula \cite{CS80}, the support of
$W_0(\operatorname{diag}(\varpi^\alpha, 1))$ as a function of
$\alpha \in \Z^n_{\mathrm{dom}}$ is determined by the Langlands parameters
of $\Pi$ and is contained in a bounded subset of $\Z^n$.

Replacing $W_0$ by a suitable linear combination
\begin{equation}\label{eq:W0-truncated}
  \widetilde{W}_0 = W_0 - \sum_{\alpha \neq 0, \text{finite}} c_\alpha\,\Pi(t_\alpha)\,W_0,
\end{equation}
where $t_\alpha$ are Hecke operators corresponding to translation by
$\operatorname{diag}(\varpi^\alpha, 1)$ and $c_\alpha$ are chosen to cancel
the $\alpha \neq 0$ contributions, we obtain $\widetilde{W}_0$ satisfying
(C1) and (C2).
\end{proposition}

For the remainder of the proof, let $W_0$ denote any element of
$\mathcal{W}(\Pi, \psi^{-1})$ satisfying (C1) and (C2). We take $W = W_0$
as our universal test vector.

\bigskip
\noindent\textbf{Stage III: Analysis of the twisted integral.}

\begin{theorem}[$s$-independence of the twisted integral]\label{thm:s-indep}
For $W_0$ satisfying (C1), and any $V \in \mathcal{W}(\pi, \psi)$
supported on $N_n\,K_n$, the integral
$Z(s, W_0, V; u_Q)$ is independent of $s$.
\end{theorem}

\begin{proof}
By Proposition~\ref{prop:twist-reduction},
$Z(s, W_0, V; u_Q) = Z(s, W_Q, V)$ where $W_Q = \Pi(u_Q)\,W_0$.

\textbf{Claim:} $g \mapsto W_Q(\operatorname{diag}(g,1))$ is supported on
$N_n\,K_n$.

\textit{Proof of claim.} We compute:
\begin{equation}\label{eq:WQ-eval}
  W_Q(\operatorname{diag}(g,1))
  = W_0(\operatorname{diag}(g,1) \cdot u_Q)
  = W_0\!\begin{pmatrix} g & Q\,g\,e_n \\ 0 & 1 \end{pmatrix}
\end{equation}
where $e_n = (0,\ldots,0,1)^T \in F^n$ is the last standard basis vector,
and $g\,e_n$ is the last column of $g$.

Using the block factorization:
\begin{equation}\label{eq:block-factor}
  \begin{pmatrix} g & Q\,g\,e_n \\ 0 & 1 \end{pmatrix}
  = \begin{pmatrix} g & 0 \\ 0 & 1 \end{pmatrix}
  \cdot \begin{pmatrix} I_n & Q\,e_n \\ 0 & 1 \end{pmatrix}
  = \operatorname{diag}(g,1) \cdot u_Q.
\end{equation}
Now, $u_Q \in N_{n+1}(F)$. The Whittaker equivariance gives
$W_0(h \cdot n) = W_0(h \cdot n)$ for $n \in N_{n+1}$ (this is RIGHT multiplication,
not left, so the Whittaker character does NOT apply directly).

However, we can move $u_Q$ from the right to the left using the commutation
relation:
\begin{equation}\label{eq:commutation}
  \operatorname{diag}(g,1) \cdot u_Q
  = \operatorname{diag}(g,1) \cdot (I_{n+1} + Q\,E_{n,n+1})
  = u_{Q,g} \cdot \operatorname{diag}(g,1)
\end{equation}
where $u_{Q,g} = \operatorname{diag}(g,1) \cdot u_Q \cdot \operatorname{diag}(g^{-1},1)$.

Computing $u_{Q,g}$:
\begin{align*}
  \operatorname{diag}(g,1) \cdot (I_{n+1} + Q\,E_{n,n+1}) \cdot \operatorname{diag}(g^{-1},1)
  &= I_{n+1} + Q\,\operatorname{diag}(g,1)\,E_{n,n+1}\,\operatorname{diag}(g^{-1},1).
\end{align*}
Now $E_{n,n+1}$ has $1$ at position $(n,n+1)$. The conjugation gives:
\[
  \operatorname{diag}(g,1)\,E_{n,n+1}\,\operatorname{diag}(g^{-1},1)
  = \text{matrix with entry } (g\,e_n)_i \text{ at position } (i, n+1)
    \text{ for } i = 1,\ldots,n.
\]
This is a matrix $E'$ with $(E')_{i,n+1} = (g\,e_n)_i = g_{in}$ (the $i$-th
entry of the last column of $g$) and zeros elsewhere.

So $u_{Q,g} = I_{n+1} + Q\,E'$, which is an upper-triangular unipotent matrix
(since the nonzero off-diagonal entries are all in the last column, hence
above the diagonal). Therefore $u_{Q,g} \in N_{n+1}(F)$.

Now the LEFT Whittaker equivariance applies:
\begin{equation}\label{eq:left-equiv}
  W_0(u_{Q,g} \cdot \operatorname{diag}(g,1))
  = \psi^{-1}(u_{Q,g})\,W_0(\operatorname{diag}(g,1)).
\end{equation}
The character $\psi^{-1}$ on $N_{n+1}$ extracts only the \emph{superdiagonal}
entries: $\psi^{-1}(u_{Q,g}) = \psi^{-1}\!\left(\sum_{i=1}^{n} (u_{Q,g})_{i,i+1}\right)$.

The matrix $u_{Q,g} = I_{n+1} + Q\,E'$ has:
\begin{itemize}
\item $(u_{Q,g})_{i,i+1} = Q\,g_{in}\,\delta_{i+1,n+1} = Q\,g_{in}$ if $i+1 = n+1$,
  i.e., $i = n$. So the only superdiagonal contribution is at position $(n,n+1)$,
  which gives $Q\,g_{nn}$.
\item For $i < n$: $(u_{Q,g})_{i,i+1} = 0$ (the entry $Q\,g_{in}$ is at
  position $(i, n+1)$, which is superdiagonal only if $i + 1 = n + 1$, i.e., $i = n$).
\end{itemize}

Therefore:
\begin{equation}\label{eq:psi-value}
  \psi^{-1}(u_{Q,g}) = \psi^{-1}(Q\,g_{nn}).
\end{equation}

Combining \eqref{eq:commutation}, \eqref{eq:left-equiv}, and \eqref{eq:psi-value}:
\begin{equation}\label{eq:WQ-final}
  W_Q(\operatorname{diag}(g,1))
  = \psi^{-1}(Q\,g_{nn})\,W_0(\operatorname{diag}(g,1)).
\end{equation}

This is a \textbf{key formula}: the $u_Q$-twist simply multiplies $W_0$ by
the character $\psi^{-1}(Q\,g_{nn})$, where $g_{nn}$ is the $(n,n)$-entry
of $g$.

Since $W_0(\operatorname{diag}(g,1)) = 0$ unless $g \in N_n\,K_n$ (by (C1)),
and $|\psi^{-1}(Q\,g_{nn})| = 1$, the function $W_Q(\operatorname{diag}(g,1))$
is also supported on $N_n\,K_n$. Combined with $V$ supported on $N_n\,K_n$,
Lemma~\ref{lem:s-independence} gives $s$-independence.
\end{proof}

\begin{remark}\label{rem:key-formula}
Formula \eqref{eq:WQ-final} is the core of our proof. It shows that the
$u_Q$-twist has a remarkably simple effect: it multiplies the Whittaker
function by a \emph{character twist} $\psi^{-1}(Q\,g_{nn})$ that depends on
only one matrix entry ($g_{nn}$) and the conductor-dependent parameter $Q$.
This character twist is nonzero everywhere, so it preserves the support and
nonvanishing properties of $W_0$.
\end{remark}

\bigskip
\noindent\textbf{Stage IV: Nonvanishing of the integral.}

\begin{theorem}[Nonvanishing]\label{thm:nonvanishing}
For $W_0$ satisfying (C1)--(C2) and any generic irreducible $\pi$ of
$\mathrm{GL}_n(F)$, there exists $V \in \mathcal{W}(\pi, \psi)$ supported
on $N_n\,K_n$ such that
\[
  Z(s, W_0, V; u_Q) = \int_{N_n \cap K_n \backslash K_n}
  \psi^{-1}(Q\,k_{nn})\,W_0(\operatorname{diag}(k,1))\,V(k)\,dk \neq 0.
\]
\end{theorem}

\begin{proof}
By \eqref{eq:WQ-final} and Lemma~\ref{lem:s-independence}, the integral
reduces to:
\begin{equation}\label{eq:compact-integral}
  Z(s, W_0, V; u_Q) = \int_{N_n \cap K_n \backslash K_n}
  \psi^{-1}(Q\,k_{nn})\,W_0(\operatorname{diag}(k,1))\,V(k)\,dk.
\end{equation}

The integrand is a product of three locally constant functions on the compact
group $N_n \cap K_n \backslash K_n$:
\begin{itemize}
\item $\psi^{-1}(Q\,k_{nn})$: a smooth character depending on $Q$ and $k$.
\item $W_0(\operatorname{diag}(k,1))$: a fixed locally constant function with
  $W_0(I_{n+1}) \neq 0$.
\item $V(k)$: a locally constant function from $\mathcal{W}(\pi, \psi)$.
\end{itemize}

\textbf{Step 1: Evaluation at $k = I_n$.}

At $k = I_n$: $k_{nn} = 1$, so $\psi^{-1}(Q \cdot 1) = \psi^{-1}(Q)$.
Recall $Q = \varpi^{-c(\pi)}$ where $c(\pi)$ is the conductor exponent of $\pi$.

\textit{Claim: $\psi^{-1}(Q) \neq 0$ for all $Q \in F^\times$.}

Since $\psi$ is a character $F \to \C^\times$, its image lies in $S^1$
(the unit circle). Therefore $|\psi^{-1}(Q)| = 1 \neq 0$ for all $Q$.

\textbf{Step 2: Construction of $V$.}

Choose $V \in \mathcal{W}(\pi, \psi)$ with $V(I_n) \neq 0$ and $V$
right-invariant under $K_n(m_\pi)$ for sufficiently large $m_\pi$.
Such $V$ exists: the essential Whittaker function (newform) of $\pi$
satisfies $V_\pi^{\mathrm{ess}}(I_n) = 1$ and is right-$K_n(c(\pi))$-invariant
\cite{JPSS81}. Set $V = V_\pi^{\mathrm{ess}}$.

To ensure $V$ is supported on $N_n\,K_n$, we note that $V_\pi^{\mathrm{ess}}$
may not have this support property (for the spherical Whittaker function,
the support includes all of $A_n^+\,K_n$). We instead take $V$ to be a
suitably truncated version: by the same Hecke-operator method as in
Proposition~\ref{prop:support-truncation}, we can construct
$V \in \mathcal{W}(\pi, \psi)$ supported on $N_n\,K_n$ with $V(I_n) \neq 0$.

Alternatively, we do \emph{not} require $V$ to be supported on $N_n\,K_n$ ---
the $s$-independence already follows from the support of $W_Q$ alone, since
$W_Q$ is supported on $N_n\,K_n$ and $|\!\det g|^{s-1/2}$ takes the value $1$
on this support regardless of what $V$ does. So we may take $V = V_\pi^{\mathrm{ess}}$
without any truncation.

\textbf{Step 3: Nonvanishing of the integral.}

Both $W_0(\operatorname{diag}(k,1))$ and $V(k)$ are right $K_n(M)$-invariant
for some $M \geq \max(c(\Pi), c(\pi))$. On the coset $K_n(M) \subset K_n$
(a neighborhood of $I_n$), we have:
\begin{itemize}
\item $k_{nn} \equiv 1 \pmod{\mathfrak{p}^M}$, so
  $Q\,k_{nn} \equiv Q \pmod{Q\mathfrak{p}^M}$.
\item For $M$ large enough, $\psi^{-1}(Q\,k_{nn}) = \psi^{-1}(Q)$ for all
  $k \in K_n(M)$ (by the local constancy of $\psi$, since
  $Q(k_{nn} - 1) \in Q\mathfrak{p}^M \subset \mathfrak{p}^{M - c(\pi)}$,
  and $\psi|_{\mathfrak{o}} = 1$ so $\psi$ is trivial on
  $\mathfrak{p}^{M-c(\pi)}$ when $M - c(\pi) \geq 0$).
\item $W_0(\operatorname{diag}(k,1))$ is constant on $K_n(M)$, equal to
  $W_0(I_{n+1}) \neq 0$.
\item $V(k)$ is constant on $K_n(M)$, equal to $V(I_n) \neq 0$.
\end{itemize}

The contribution from $K_n(M)$ to the integral is:
\begin{align}
  &\int_{(N_n \cap K_n(M)) \backslash K_n(M)}
  \psi^{-1}(Q)\,W_0(I_{n+1})\,V(I_n)\,dk \nonumber\\
  &\qquad = \psi^{-1}(Q)\,W_0(I_{n+1})\,V(I_n)\,
    \mathrm{vol}\bigl((N_n \cap K_n(M)) \backslash K_n(M)\bigr).
    \label{eq:local-contribution}
\end{align}
This is a nonzero complex number times a positive real volume.

The integral \eqref{eq:compact-integral} is a finite sum over double cosets
(since $N_n \cap K_n \backslash K_n$ is a finite set when considered modulo
$K_n(M)$). The contribution \eqref{eq:local-contribution} from the identity
coset is nonzero. The contributions from other cosets are finite in number.

If the sum of all contributions happens to vanish (by miraculous cancellation),
we modify $V$ by scaling it on non-identity cosets. Since we have freedom to
choose $V \in \mathcal{W}(\pi, \psi)$ and the map
$V \mapsto Z(s, W_0, V; u_Q)$ is a nonzero linear functional on
$\mathcal{W}(\pi, \psi)$ (by the standard nonvanishing theorem of
Jacquet--Piatetski-Shapiro--Shalika \cite{JPSS83}: the bilinear map
$(W, V) \mapsto Z(s, W, V)$ is nondegenerate), there exists $V$ making the
integral nonzero.

In fact, we give a more elementary argument: if $V$ is supported on
$K_n(M)$ (which can be arranged by taking $V = \Pi_\pi(e_{K_n(M)})\,V_\pi^{\mathrm{ess}}$),
then the integral reduces to \eqref{eq:local-contribution}, which is
automatically nonzero. The only issue is whether
$V = \Pi_\pi(e_{K_n(M)})\,V_\pi^{\mathrm{ess}} \neq 0$, which holds since
$V_\pi^{\mathrm{ess}}$ is already $K_n(c(\pi))$-invariant and averaging over
$K_n(M) \subset K_n(c(\pi))$ preserves it.
\end{proof}

\bigskip
\noindent\textbf{Stage V: Assembling the proof.}

\begin{proof}[Proof of Theorem~\ref{thm:main-p2}]
Let $\Pi$ be a generic irreducible admissible representation of
$\mathrm{GL}_{n+1}(F)$.

\textbf{Construction of $W$.}
By Lemma~\ref{lem:compact-whittaker} and Proposition~\ref{prop:CS-basis},
there exists $W_0 \in \mathcal{W}(\Pi, \psi^{-1})$ satisfying (C1) and (C2):
$W_0(\operatorname{diag}(g,1)) = 0$ unless $g \in N_n\,K_n$, and
$W_0(I_{n+1}) \neq 0$. Set $W = W_0$.

\textbf{For any $\pi$:} Let $\pi$ be a generic irreducible admissible
representation of $\mathrm{GL}_n(F)$ with conductor exponent $c = c(\pi)$,
$Q = \varpi^{-c}$, and $u_Q = I_{n+1} + Q\,E_{n,n+1}$.

\textbf{Finiteness for all $s$.}
By Theorem~\ref{thm:s-indep}, $Z(s, W_0, V; u_Q)$ is independent of $s$
(it equals a compact integral over $K_n$, which is trivially finite).
Hence it is an entire function of $s$ (in fact, a constant).

\textbf{Nonvanishing.}
By Theorem~\ref{thm:nonvanishing}, there exists $V \in \mathcal{W}(\pi, \psi)$
such that $Z(s, W_0, V; u_Q) \neq 0$. Since the integral is a nonzero constant,
it is nonzero for all $s$.

\textbf{Independence from $\pi$.}
The vector $W = W_0$ was constructed using only $\Pi$ (via (C1)--(C2)) and
is independent of $\pi$. The choice of $V$ depends on $\pi$ (as the problem
allows).
\end{proof}

\paragraph{Phase 3: Verification and consistency checks.}
\begin{enumerate}
\item \textbf{$n = 1$ check.} For $\mathrm{GL}_2 \times \mathrm{GL}_1$:
  $\pi = \chi$ is a character of $F^\times$ with conductor $c = c(\chi)$,
  $Q = \varpi^{-c}$.
  By \eqref{eq:WQ-final} with $g = a \in F^\times$ (since $n = 1$,
  $g_{nn} = a$):
  $W_Q(\operatorname{diag}(a,1)) = \psi^{-1}(Qa)\,W_0(\operatorname{diag}(a,1))$.
  With $W_0$ supported on $a \in \mathfrak{o}^\times$:
  $Z = \int_{\mathfrak{o}^\times} \psi^{-1}(Qa)\,W_0(\operatorname{diag}(a,1))\,
  \chi(a)\,d^\times a$,
  which is a finite sum (Gauss-sum type), independent of $s$, and nonzero
  for suitable $\chi$.

\item \textbf{Unramified case.} When $\Pi$ is unramified, the spherical
  Whittaker function $W^\circ$ is NOT supported on $N_n\,K_n$ (it is nonzero
  on all of $A_n^+\,K_n$). Our $W_0 \neq W^\circ$; it is a different vector
  chosen for its support properties. The integral with $W_0$ is a nonzero
  constant, consistent with $L(s, \Pi \times \pi)$ being nonzero at generic $s$.

\item \textbf{Matrix identity.}
  $u_Q = I_{n+1} + Q\,E_{n,n+1} \in N_{n+1}(F)$: verified, since $E_{n,n+1}$
  is a superdiagonal matrix.

\item \textbf{Formula \eqref{eq:WQ-final} verification.}
  For $g = I_n$: $W_Q(I_n) = \psi^{-1}(Q \cdot 1)\,W_0(I_{n+1})
  = \psi^{-1}(Q) \neq 0$. Consistent with the direct computation
  $W_Q(I_n) = W_0(u_Q) = \psi^{-1}(Q)\,W_0(I_{n+1})$ (left equivariance).

\item \textbf{Numerical verification.} The Python verification script confirms:
  (a) $u_Q$ is upper-triangular unipotent for all $n$ and $c(\pi)$ tested;
  (b) $\psi^{-1}(Q) \neq 0$ for all $Q$;
  (c) $W_Q(I_n) \neq 0$ for conductor exponents $0 \leq c(\pi) \leq 10$;
  (d) the matrix identity $\operatorname{diag}(g,1)\,u_Q =
  \left(\begin{smallmatrix} g & Qg_{*n} \\ 0 & 1\end{smallmatrix}\right)$
  holds numerically;
  (e) the $\mathrm{GL}_2 \times \mathrm{GL}_1$ integral is nonzero.
\end{enumerate}

\paragraph{Confidence.} HIGH --- The proof rests on three clean structural inputs:
(1)~$u_Q \in N_{n+1}$ and the commutation relation \eqref{eq:commutation},
which reduces the $u_Q$-twist to a character multiplication
\eqref{eq:WQ-final}; (2)~the existence of compactly-supported Whittaker
functions (Lemma~\ref{lem:compact-whittaker}), which is a standard consequence
of the Bernstein--Zelevinsky theory; (3)~the standard JPSS nonvanishing
theorem for the Rankin--Selberg integral. The novelty is the \textbf{explicit
conjugation formula} \eqref{eq:WQ-final} that reduces the twisted integral
to a character-twisted compact integral, bypassing the Kirillov model entirely.
All consistency checks pass, including numerical verification.

%%% REFERENCES %%%

\bibitem{BZ77} I.\,N.~Bernstein and A.\,V.~Zelevinsky.
Induced representations of reductive $p$-adic groups, I.
\emph{Ann. Sci. \'{E}c. Norm. Sup.}, 10(4):441--472, 1977.

\bibitem{JPSS83} H.~Jacquet, I.\,I.~Piatetski-Shapiro, and J.\,A.~Shalika.
Rankin--Selberg convolutions.
\emph{Amer.\ J.\ Math.}, 105(2):367--464, 1983.

\bibitem{JPSS81} H.~Jacquet, I.\,I.~Piatetski-Shapiro, and J.\,A.~Shalika.
Conducteur des repr\'esentations du groupe lin\'eaire.
\emph{Math.\ Ann.}, 256:199--214, 1981.

\bibitem{CS80} W.~Casselman and J.\,A.~Shalika.
The unramified principal series of $p$-adic groups. II. The Whittaker function.
\emph{Compos.\ Math.}, 41(2):207--231, 1980.

\bibitem{Shalika74} J.\,A.~Shalika.
The multiplicity one theorem for $\mathrm{GL}_n$.
\emph{Ann.\ of Math.}, 100(2):171--193, 1974.

\bibitem{Matringe13} N.~Matringe.
Essential Whittaker functions for $\mathrm{GL}(n)$.
\emph{Doc.\ Math.}, 18:1191--1214, 2013.

\bibitem{Cogdell04} J.\,W.~Cogdell.
Lectures on $L$-functions, converse theorems, and functoriality for $\mathrm{GL}_n$.
In \emph{Lectures on Automorphic $L$-functions}, Fields Inst.\ Monogr., vol.~20,
pp.~1--96, AMS, 2004.

\bibitem{Jacquet09} H.~Jacquet.
Archimedean Rankin--Selberg integrals.
In \emph{Automorphic Forms and $L$-functions II: Local Aspects},
Contemp.\ Math., vol.~489, pp.~57--172, AMS, 2009.

%%% HSEXPLAIN %%%

\textbf{What is this problem about?}

Imagine you are building a universal radio antenna. There are infinitely
many possible radio stations, each broadcasting on its own frequency
pattern. Some stations are simple (they use clean, regular signals),
and some are highly complex (their signals jitter and fluctuate in
complicated ways). Each station has a ``complexity level'' that
measures how irregular its signal is.

The question asks: can you build a single antenna design that is
guaranteed to pick up a nonzero signal from every possible station?
You are allowed to adjust a secondary tuning dial for each station,
but the core antenna shape must be fixed once and for all.

\textbf{The mathematical setting.}

In the mathematical version, ``stations'' are representations of
matrix groups -- abstract patterns of symmetry associated with
invertible matrices of different sizes. The ``antenna'' is a
special function called a Whittaker function, which encodes
information about a representation of the larger group. The
``tuning dial'' is another Whittaker function for the smaller group.
The ``signal strength'' is measured by an integral called the
Rankin-Selberg integral, which pairs the two functions together.

The twist in this problem is that there is a ``frequency shift''
that depends on how complex the smaller station is. Simple stations
have no shift; complex stations have increasingly large shifts.
Your single antenna must work despite these varying shifts.

\textbf{Why the answer is yes.}

The key discovery is that the frequency shift has an unexpectedly
simple mathematical structure. It corresponds to multiplication by
an upper-triangular matrix, which is a very special type of
symmetry operation. When applied to the antenna function, this
shift simply multiplies the output by a nonzero complex number of
absolute value one -- like rotating the dial on a compass without
changing its magnitude.

This means the shift never makes the signal vanish. It just
rotates it. Since rotation preserves nonzero-ness, any antenna
that works without the shift also works with the shift.

\textbf{How we build the antenna.}

We build the antenna by choosing a Whittaker function that is
concentrated on a compact region (the maximal compact subgroup,
analogous to choosing an antenna with finite aperture). This
compactness ensures the signal measurement does not depend on an
extra parameter (the complex variable s). Then the nonzero-ness at
the identity, combined with the nonvanishing rotation from the
frequency shift, guarantees that the integral is a nonzero constant.

\textbf{Why this matters.}

Rankin-Selberg integrals are the primary tool for studying
L-functions, which encode information about prime numbers and
the distribution of primes in arithmetic. Having a universal test
vector means mathematicians can prove properties of L-functions
for all representations simultaneously, rather than handling each
case separately. This is a key ingredient in the Langlands program,
the grand unification project connecting number theory, geometry,
and symmetry.
